\documentclass[11pt]{book}

% PRIMARY ENGINE: lualatex main.tex (two passes mandatory for cross-references)
% FALLBACK ENGINE: xelatex main.tex (if LuaLaTeX microtype conflict arises)
%
% Known LuaLaTeX issue: microtype may conflict with TeX Gyre Pagella under
% specific font-backend configurations. If LuaLaTeX fails, comment out
% \usepackage{microtype} in preamble.tex and retry, or use XeLaTeX fallback.
% Both engines supported via fontspec; no engine-specific features used.

% preamble.tex
% Shared preamble for "Admissible Reconstruction"
% Engine-agnostic via fontspec (compiles under LuaLaTeX or XeLaTeX)

\usepackage{fontspec}
\setmainfont{TeX Gyre Pagella}

\usepackage{amsmath}
\usepackage{amssymb}
\usepackage{amsthm}
\usepackage{mathtools}
\usepackage[margin=1.25in]{geometry}
\usepackage[colorlinks=true,linkcolor=black,citecolor=black,urlcolor=black]{hyperref}
\usepackage{microtype}
\usepackage{enumitem}
\usepackage{titlesec}

% ----- Theorem environments -----
\theoremstyle{plain}
\newtheorem{theorem}{Theorem}[chapter]
\newtheorem{proposition}[theorem]{Proposition}
\newtheorem{lemma}[theorem]{Lemma}
\newtheorem{corollary}[theorem]{Corollary}

\theoremstyle{definition}
\newtheorem{definition}[theorem]{Definition}
\newtheorem{example}[theorem]{Example}

\theoremstyle{remark}
\newtheorem{remark}[theorem]{Remark}

% ----- Notation macros -----
% Admissibility relation x ~>_R y
\newcommand{\adm}[1]{\rightsquigarrow_{#1}}
% Reachable set of a state
\newcommand{\reach}[1]{R(#1)}
% Time-indexed admissibility structure
\newcommand{\Rt}{R_t}
% Closure functional
\newcommand{\closure}[1]{C(#1)}
% Generativity measure
\newcommand{\genmeas}[1]{G(#1)}
% Observation map
\newcommand{\obsmap}{O}
% History / observable / reconstruction
\newcommand{\History}{H}
\newcommand{\Observable}{O_t}
\newcommand{\Reconstruction}{\hat{H}}

% ----- Section numbering depth -----
\setcounter{secnumdepth}{3}
\setcounter{tocdepth}{2}

% ----- No bulleted lists in body prose: enumitem loaded for the rare
% cases (definitions, notation index) where a list is structurally
% necessary rather than a substitute for prose. -----


\title{\textbf{Admissible Reconstruction}}
\author{Flyxion \\ \textit{Independent Researcher}}
\date{}

\begin{document}

\frontmatter

\maketitle

\chapter*{Abstract}
\addcontentsline{toc}{chapter}{Abstract}

This monograph isolates and proves four theorems governing when structure
survives transformation, and when it does not. Each theorem is derived
directly from the admissibility apparatus developed elsewhere in this
research program --- the relation $x \adm{R} y$, the reachable set
$\reach{x}$, the time-indexed admissibility structure $R_t$, and the
closure functional $\closure{X} = -\log \mu(\reach{X})$ --- rather than
introduced as free-standing results. \emph{Representation Invariance}
states that two representations are semantically equivalent whenever
they induce the same admissibility structure, independent of their
physical substrate. \emph{Admissibility Preservation} shows that a
trajectory's coherence is not an independent postulate but a consequence
of the admissibility relation failing to degenerate: coherence just is
the nontriviality of $R(S_t)$ maintained across a trajectory.
\emph{Monotonic Relaxation} identifies the precise, conditional sense in
which a scalar measure of inadmissibility decreases under repair
dynamics specifically, and distinguishes this claim carefully from the
unrelated ordering relation of monotonic continuation. \emph{Non-Identifiability}
shows that the existence of histories $H_1 \neq H_2$ with
$\obsmap(H_1) = \obsmap(H_2)$ is not a pathological edge case but a
structural consequence of any observation map that preserves only
reachable-set structure rather than full state --- the same fact that
elsewhere in this program is called the restorability boundary, here
proven as a limitative theorem in its own right.

Together the four results describe a single underlying architecture:
representation, coherence, repair, and the limits of recovery are not
four separate phenomena requiring four separate theories, but four
faces of one admissibility relation examined at different points in its
behavior.

\tableofcontents

\chapter*{Notation Index}
\addcontentsline{toc}{chapter}{Notation Index}

\begin{description}[leftmargin=2.2cm, style=nextline]
  \item[$x \adm{R} y$] $x$ admissibly continues to $y$ under admissibility
    structure $R$.
  \item[$\reach{x}$] the reachable set of $x$: $\{\, y : x \adm{R} y \,\}$.
  \item[$R_t$] the admissibility structure indexed by time $t$; the
    static relation $R$ is the frozen special case $R_t \equiv R$.
  \item[$S_t \le S_{t+1}$] monotonic continuation: $S_{t+1}$ preserves
    the distinctions of $S_t$ that remain necessary for the reachable
    trajectory. Ordering, not magnitude.
  \item[$\closure{X} = -\log \mu(\reach{X})$] the closure functional; a
    scalar measure of inadmissibility, larger when fewer futures remain
    reachable from $X$.
  \item[$\genmeas{X} = \log \mu(\reach{X})$] the generativity measure;
    $\genmeas{X} = -\closure{X}$.
  \item[$\History$] a complete history (ground truth), as retained by a
    simulator or, more generally, as the full state trajectory prior to
    any lossy observation.
  \item[$\Observable$] the observable projection of $\History$ available
    at time $t$; a deliberately lossy map $\obsmap : \History \to
    \Observable$.
  \item[$\Reconstruction$] a reconstruction of $\History$ inferred from
    $\Observable$ alone, without privileged access to $\History$.
  \item[$\varphi$] an admissibility-isomorphism between two
    representations: a map such that $x \adm{R} y \iff \varphi(x)
    \adm{R'} \varphi(y)$.
\end{description}

\mainmatter

\chapter{Introduction}
\label{ch:introduction}

This monograph is organized around a single question: when does structure
survive transformation, and when does it not? The question has appeared
throughout this research program under several names --- reconstruction,
repair, restorability, continuation --- and has been approached from
several directions, including a formal treatment of admissibility and
irreversibility, a foundational account of observability as restorability,
and an extended study of continuation across unequal timescales. What
follows is not a new direction but a consolidation: four theorems, each
of which has appeared informally or partially elsewhere in this program,
stated here as precise results and derived directly from the admissibility
apparatus rather than argued for case by case.

The apparatus itself is simple to state. An admissibility structure $R$
determines, for any state $x$, a reachable set $\reach{x}$ of states $y$
such that $x \adm{R} y$. Nothing about $R$ is assumed beyond this: it is
a relation, not a metric, and it need not be symmetric, transitive, or
total. Historical work in this program has shown that a great deal of
structure can be recovered from this minimal starting point once it is
allowed to vary with time, $R_t$, and once transitions between admissible
structures are themselves treated as admissible or inadmissible events
rather than as external clock ticks.

The four theorems developed in this monograph correspond to four
questions one can ask of this apparatus. First: when are two
representations, possibly built from entirely different physical or
computational substrates, the same object as far as $R$ is concerned?
This is Representation Invariance, the subject of Chapter~\ref{ch:invariance}.
Second: what does it mean for a trajectory through admissibility structure
to be coherent, and is coherence an additional postulate or a consequence
of the relation itself? This is Admissibility Preservation, the subject
of Chapter~\ref{ch:preservation}. Third: under what conditions, and in
what precise sense, does a scalar measure of inadmissibility decrease
along a trajectory, and how does this relate to --- and how must it be
distinguished from --- the ordering relation of monotonic continuation
developed elsewhere in this program? This is Monotonic Relaxation, the
subject of Chapter~\ref{ch:relaxation}. Fourth: when is it structurally
impossible, not merely difficult, to recover which of two histories
actually occurred from the evidence either would have left behind? This
is Non-Identifiability, the subject of Chapter~\ref{ch:nonidentifiability}.

A methodological note is appropriate before proceeding. Each of the four
theorems could be, and in earlier work occasionally has been, treated as
an independent claim requiring its own justification. The approach taken
here is different. Every theorem below is derived from definitions
already established elsewhere in this program --- the admissibility
relation, the reachable set, the closure functional --- so that what is
new in this monograph is not the vocabulary but the demonstration that
four apparently distinct phenomena are, in fact, four views of one
relation. Where a theorem's content is close to definitional, this is
stated plainly rather than disguised as a deeper result; where a theorem
requires a genuine argument, as with Admissibility Preservation and
Monotonic Relaxation, the argument is given in full and its conditions
of validity are stated explicitly rather than left implicit.

\chapter{Representation Invariance}
\label{ch:invariance}

\section{Statement}

The question addressed in this chapter is when two representations,
possibly realized in entirely different substrates, ought to be
considered the same object. The answer given here is that sameness of
substrate is irrelevant; what matters is sameness of admissibility
structure.

\begin{definition}[Admissibility-isomorphism]
Let $R$ be an admissibility structure on a set $X$ and $R'$ an
admissibility structure on a set $X'$. A map $\varphi : X \to X'$ is an
\emph{admissibility-isomorphism} if, for all $x, y \in X$,
\[
x \adm{R} y \quad \iff \quad \varphi(x) \adm{R'} \varphi(y).
\]
\end{definition}

\begin{theorem}[Representation Invariance]
\label{thm:invariance}
If there exists an admissibility-isomorphism $\varphi : X \to X'$
between $(X, R)$ and $(X', R')$, then $(X, R)$ and $(X', R')$ are
semantically equivalent, regardless of whatever physical or
computational substrate realizes $X$ and $X'$.
\end{theorem}

\section{Discussion}

The theorem as stated is close to definitional, and this is by design
rather than by accident. Semantic equivalence, within this research
program, has never been defined as identity of physical encoding; it has
been defined, throughout the treatment of admissible representation and
context preservation developed elsewhere in this program, as identity of
what a representation permits and forbids as continuation. Theorem
\ref{thm:invariance} simply makes that definition explicit as an
invariance statement rather than leaving it distributed across the
apparatus that depends on it.

What is not definitional, and what does real work when the theorem is
applied, is establishing that a particular $\varphi$ is in fact an
admissibility-isomorphism between two representations of interest. This
is generally the harder direction: showing that a biphasic transformation
--- structured meaning realized as continuous speech, or as sequential
text, or as embodied gesture --- preserves the reachable-set structure
exactly, rather than merely preserving some coarser statistic that
happens to look similar under casual inspection. The content of a
representation-invariance claim, in any particular application, is
carried entirely by the verification that $\varphi$ satisfies the
biconditional in both directions: that admissible continuations in one
representation correspond exactly to admissible continuations in the
other, with neither representation admitting a continuation the other
forbids.

\begin{remark}
It follows immediately from Theorem \ref{thm:invariance} that
representation invariance is an equivalence relation on admissibility
structures: the identity map witnesses reflexivity, the inverse of an
admissibility-isomorphism is again an admissibility-isomorphism, giving
symmetry, and composition of admissibility-isomorphisms is again an
admissibility-isomorphism, giving transitivity. Consequently, the class
of representations equivalent to a given $(X, R)$ partitions into a
single equivalence class under $\varphi$, and asking whether two
representations are "the same" reduces to asking whether they lie in the
same class.
\end{remark}

\begin{corollary}
\label{cor:invariance-substrate}
No property of $X$ or $X'$ that is not preserved by every
admissibility-isomorphism can be a semantic property of the
representation. In particular, substrate-specific properties --- the
particular physical channel through which a representation is realized,
its encoding cost, or its production latency --- are not semantic
properties under this framework unless they happen to be expressible as
constraints on $R$ itself.
\end{corollary}

Corollary \ref{cor:invariance-substrate} is worth stating explicitly
because it draws a sharp boundary that is easy to blur in practice.
Structural semantics, as developed elsewhere in this program, is
precisely the study of what survives this partition; anything that does
not survive it belongs to a different, and legitimate, subject ---
the study of channel, cost, and implementation --- but not to the study
of what a representation \emph{means}.

\chapter{Admissibility Preservation}
\label{ch:preservation}

\section{Statement}

The question addressed in this chapter is what it means for a
trajectory through admissibility structure to remain coherent. The
answer given here is that coherence is not an additional condition
imposed on top of admissibility; it is what nontriviality of the
admissibility relation amounts to when examined along a trajectory
rather than at a single state.

\begin{definition}[Trajectory]
A trajectory is a sequence of states $S_0, S_1, \dots, S_t$ such that
$S_i \adm{R} S_{i+1}$ for each $i < t$, where $R$ may itself vary with
$i$ if the underlying admissibility structure is time-indexed as $R_i$.
\end{definition}

\begin{definition}[Degeneracy]
A trajectory degenerates at step $i$ if $\reach{S_i}$ is empty (no
continuation is admissible) or a singleton entirely determined by
$S_0, \dots, S_{i-1}$ (exactly one continuation is admissible, and it
was already fixed before $S_i$ was reached).
\end{definition}

\begin{theorem}[Admissibility Preservation]
\label{thm:preservation}
A trajectory $S_0, \dots, S_t$ is coherent if and only if it does not
degenerate at any step $i \le t$; equivalently, if and only if
$|\reach{S_i}|$ exceeds a nontrivial floor relative to what has already
been committed by $S_0, \dots, S_{i-1}$, at every step.
\end{theorem}

\begin{proof}
Coherence, as the term is used throughout this program, has no content
apart from the admissibility relation itself: to say a trajectory is
coherent is to say that at each step, the relation $x \adm{R} y$
continues to discriminate admissible from inadmissible continuations
non-vacuously. Suppose the trajectory degenerates at step $i$ because
$\reach{S_i} = \varnothing$. Then no $y$ satisfies $S_i \adm{R} y$, and
the relation, restricted to continuations from $S_i$, is uniformly
false; there is no sense in which the trajectory can be said to continue
admissibly past $S_i$, coherently or otherwise, since no continuation is
admissible at all. Suppose instead the trajectory degenerates because
$\reach{S_i}$ is a singleton already fixed by $S_0, \dots, S_{i-1}$.
Then the relation restricted to continuations from $S_i$ is uniformly
true of exactly one candidate and uniformly false of every other; the
admissibility relation is doing no discriminating work at step $i$, since
its outcome was determined before $S_i$ was reached, and the statement
"$S_i \adm{R} S_{i+1}$ is coherent" becomes true independently of
anything about $S_i$ itself. In either degenerate case, the relation
$\adm{R}$ fails to be a nontrivial constraint at step $i$: it is either
vacuously unsatisfiable or vacuously predetermined. Conversely, if
$\reach{S_i}$ exceeds a nontrivial floor at every step --- if there is
always more than one admissible continuation not already fixed by the
trajectory's own history --- then the admissibility relation is doing
genuine discriminating work at every step, which is what coherence
consists in. This establishes the biconditional.
\end{proof}

\section{Discussion}

The proof above is deliberately minimal, and its minimality is the
point. Coherence has sometimes been treated, informally, as though it
were an additional property a trajectory could have or lack independent
of the admissibility relation that generates it --- as though a
trajectory could be admissible at every step and yet somehow
incoherent, or coherent despite violating admissibility. Theorem
\ref{thm:preservation} rules this out: coherence and non-degeneracy of
the admissibility relation are the same property under two names. This
matters directly for the definition of monotonic continuation developed
elsewhere in this program, where the ordering $S_t \le S_{t+1}$
presupposes, without stating, that $\reach{S_t}$ remains nontrivial;
Theorem \ref{thm:preservation} supplies the missing lemma that makes
this presupposition explicit and shows exactly what fails when it does
not hold.

\begin{corollary}
\label{cor:preservation-minimum-support}
Any operational test for coherence --- such as a minimum support ratio
required of a discourse trajectory before it is judged coherent --- is
correctly understood as an operational proxy for $|\reach{S_i}|$
remaining above a nontrivial floor, not as an independent criterion.
Where such a proxy diverges from the underlying reachable-set condition,
the proxy, not the theorem, should be revised.
\end{corollary}

Corollary \ref{cor:preservation-minimum-support} is included because
domain-specific coherence tests --- support ratios, consistency scores,
and similar quantities --- are frequently introduced as though they were
primitive definitions of coherence in their own right. Theorem
\ref{thm:preservation} reduces the space of legitimate such tests to
those that track, however indirectly, the nondegeneracy of the
admissibility relation, and gives a principled way to evaluate a
proposed test: ask whether it can be false while $\reach{S_i}$ remains
nontrivial, or true while $\reach{S_i}$ has degenerated. Either case
indicates the proxy has diverged from what coherence actually
consists in.

\chapter{Monotonic Relaxation}
\label{ch:relaxation}

\section{A Terminological Warning}

This chapter's theorem shares a name, and very little else, with
monotonic continuation as developed elsewhere in this program. The two
should not be conflated, and the distinction is stated here before the
theorem itself to prevent exactly that conflation. Monotonic
continuation is an ordering relation, $S_t \le S_{t+1}$, holding between
successive states of a trajectory and asserting that later states
preserve the distinctions from earlier states that remain necessary for
the reachable continuation; it says nothing about any scalar quantity
increasing or decreasing. Monotonic relaxation, the subject of this
chapter, is a claim that a specific scalar functional decreases along a
specific class of trajectories. The two claims are compatible, and one
will be shown below to imply a restricted version of the other under an
additional hypothesis, but they are not restatements of each other and
a trajectory can satisfy either without satisfying the other.

\section{Statement}

\begin{definition}[Closure functional]
For a state $X$, define $\closure{X} = -\log \mu(\reach{X})$, where
$\mu$ is a measure on the state space assigning nonzero mass to any
nonempty reachable set. $\closure{X}$ is large when few futures remain
reachable from $X$ and small when many do; it is a scalar measure of
inadmissibility.
\end{definition}

\begin{definition}[Repair dynamics]
A trajectory follows \emph{repair dynamics} if each transition
$S_i \to S_{i+1}$ is selected, among admissible continuations, to
increase reachable-set volume: $\mu(\reach{S_{i+1}}) \ge
\mu(\reach{S_i})$. This is the operational content of repair within this
program's taxonomy of admissibility-modifying operations: repair is
the restorative case, distinguished from constitutive and regulatory
modification precisely by this volume-increasing character.
\end{definition}

\begin{theorem}[Monotonic Relaxation]
\label{thm:relaxation}
If a trajectory $S_0, \dots, S_t$ follows repair dynamics throughout,
then $\closure{S_i}$ is non-increasing along the trajectory:
$\closure{S_{i+1}} \le \closure{S_i}$ for every $i < t$.
\end{theorem}

\begin{proof}
By definition of repair dynamics, $\mu(\reach{S_{i+1}}) \ge
\mu(\reach{S_i})$ for every $i < t$. Since $-\log$ is a strictly
decreasing function on positive reals, $\mu(\reach{S_{i+1}}) \ge
\mu(\reach{S_i})$ implies $-\log \mu(\reach{S_{i+1}}) \le -\log
\mu(\reach{S_i})$, that is, $\closure{S_{i+1}} \le \closure{S_i}$.
\end{proof}

\section{Discussion}

The proof is immediate once repair dynamics are defined as
volume-increasing, and this is deliberate: the theorem's content is not
in the inference from the definition, which is a one-line consequence of
$-\log$ being decreasing, but in the definition itself being the correct
operational characterization of repair. The definition is not
stipulated arbitrarily; it is inherited directly from the restorative
case of this program's modificatory-operation taxonomy, where repair is
already characterized, independently of this chapter, as the operation
type that restores lost admissible continuations. Theorem
\ref{thm:relaxation} shows that once repair is characterized this way,
monotone decrease of $\closure{\cdot}$ follows automatically; it is not
an additional empirical claim about repair processes but a restatement
of what "restoring lost continuations" already meant, now expressed as
a statement about a scalar functional.

\begin{remark}[Scope of the theorem]
The hypothesis is essential and is not automatically satisfied by an
arbitrary admissible trajectory. Constitutive and regulatory
modifications --- the other two branches of this program's taxonomy of
operations on admissibility structure --- need not increase
$\mu(\reach{\cdot})$, and ordinary state evolution under a fixed $R_t$
can move toward states with strictly smaller reachable sets without
this constituting any failure of admissibility. A claim of the form
"inadmissibility always decreases" would be false; Theorem
\ref{thm:relaxation} is true only of the restorative subclass of
trajectories, and its scope should be stated with this restriction every
time it is invoked.
\end{remark}

\begin{corollary}[Conditional relation to monotonic continuation]
\label{cor:relaxation-continuation}
If a trajectory follows repair dynamics and additionally satisfies
$\reach{S_{i+1}} \supseteq^{*} \reach{S_i}$ in the starred sense used by
monotonic continuation --- that is, $S_{i+1}$ retains the admissible
continuations of $S_i$ that remain significant for the reachable
trajectory, rather than merely having a reachable set of at least equal
measure --- then the trajectory satisfies both $S_i \le S_{i+1}$ in the
sense of monotonic continuation and $\closure{S_{i+1}} \le \closure{S_i}$
in the sense of this chapter. Neither condition implies the other in
general: a trajectory can increase $\mu(\reach{\cdot})$ while failing to
retain the specific significant continuations monotonic continuation
requires, by replacing them with an equal or greater volume of
differently located reachable states; and a trajectory can satisfy
monotonic continuation's retention condition while $\mu(\reach{\cdot})$
decreases, if what is retained is a shrinking but still significant
core.
\end{corollary}

Corollary \ref{cor:relaxation-continuation} is the precise statement of
the relationship gestured at in the opening warning of this chapter.
The two notions of monotonicity coincide only under an additional
hypothesis that must be checked independently in any given application,
and the corollary exists chiefly to prevent the two theorems from being
cited interchangeably.

\chapter{Non-Identifiability}
\label{ch:nonidentifiability}

\section{Statement}

The question addressed in this chapter is when it is impossible, not
merely difficult, to recover which of two histories actually occurred
from the evidence either would have left behind. The elementary
direction of this claim is close to definitional; the substantive
direction is showing that the relevant class of observation maps is
genuinely, structurally non-injective, rather than merely
non-injective in unusual or contrived cases.

\begin{definition}[Observation map]
An observation map is a function $\obsmap : \History \to \Observable$
from complete histories to observables, where $\Observable$ is whatever
evidence remains available at a given time to an observer without
access to $\History$ itself.
\end{definition}

\begin{theorem}[Non-Identifiability, elementary direction]
\label{thm:nonident-elementary}
If $H_1 \neq H_2$ and $\obsmap(H_1) = \obsmap(H_2)$, then no function
$f : \Observable \to \History$ satisfies $f(\obsmap(H_1)) = H_1$ and
$f(\obsmap(H_2)) = H_2$ simultaneously.
\end{theorem}

\begin{proof}
Suppose such an $f$ existed. Since $\obsmap(H_1) = \obsmap(H_2)$, write
$o = \obsmap(H_1) = \obsmap(H_2)$. Then $f(o) = f(\obsmap(H_1)) = H_1$
and also $f(o) = f(\obsmap(H_2)) = H_2$, so $H_1 = H_2$, contradicting
the hypothesis $H_1 \neq H_2$.
\end{proof}

This direction is elementary because it is a direct consequence of $f$
being required to be a function, together with the assumed collision
$\obsmap(H_1) = \obsmap(H_2)$; no property of admissibility structure is
needed to prove it. What requires the admissibility apparatus developed
elsewhere in this program is establishing that such collisions are not
pathological exceptions confined to contrived cases, but a structural
feature of the observation maps that actually arise when observation is
constrained to preserve only reachable-set structure.

\begin{theorem}[Non-Identifiability, structural direction]
\label{thm:nonident-structural}
Let $\obsmap$ be any observation map that factors through the
reachable-set structure --- that is, suppose $\obsmap(H) = \psi(\reach{H})$
for some function $\psi$, so that $\obsmap$ records only which
continuations remain admissible from $H$, not $H$ itself. Then whenever
two histories $H_1 \neq H_2$ satisfy $\reach{H_1} = \reach{H_2}$, it
follows that $\obsmap(H_1) = \obsmap(H_2)$, and Theorem
\ref{thm:nonident-elementary} applies. Moreover, this is not a
degenerate case: $\reach{H_1} = \reach{H_2}$ is precisely the condition
under which $H_1$ and $H_2$ lie in the same equivalence class of the
admissibility relation, and any admissibility structure with a
non-trivial quotient --- that is, any admissibility structure in which
more than one history maps to the same reachable set --- guarantees the
existence of such collisions.
\end{theorem}

\begin{proof}
The first claim is immediate: if $\reach{H_1} = \reach{H_2}$ then
$\psi(\reach{H_1}) = \psi(\reach{H_2})$ trivially, since $\psi$ is a
function of the reachable set alone. For the second claim, recall that
$R(\cdot)$ partitions the space of histories into equivalence classes by
definition: two histories lie in the same class exactly when they have
the same reachable set. An admissibility structure has a trivial
quotient only when this partition is the discrete partition, that is,
only when no two distinct histories ever share a reachable set. This is
not the generic case for admissibility structures of the kind developed
elsewhere in this program, where reachable sets are typically determined
by a bounded amount of structure --- the admissible continuations
consistent with what has been committed so far --- rather than by the
full, unbounded historical record. Whenever the reachable set is a
function of bounded information about $H$ while $H$ itself can vary over
an unbounded space, the pigeonhole principle guarantees histories
differing in the unrecorded portion while agreeing in the recorded
portion, hence sharing a reachable set, hence colliding under $\obsmap$.
\end{proof}

\section{Discussion}

Theorem \ref{thm:nonident-structural} is the precise sense in which
Non-Identifiability is not a separate result requiring separate
machinery, but the restorability boundary developed elsewhere in this
program, proven here as a limitative theorem in its own right rather
than argued for informally. The restorability boundary has previously
been discussed in terms of what invariants survive repeated
reconstruction and what information is irrecoverably lost; Theorem
\ref{thm:nonident-structural} identifies exactly the mechanism by which
loss becomes structurally guaranteed rather than merely likely: any
observation map constrained to record reachable-set structure, applied
to a space of histories whose reachable sets are determined by less
information than the histories themselves contain, is non-injective as
a matter of definition, not as a matter of contingent bad luck in any
particular case.

\begin{corollary}
\label{cor:nonident-diagnosis}
Given an observed collision $\obsmap(H_1) = \obsmap(H_2)$ with
$H_1 \neq H_2$, the question of \emph{why} reconstruction fails for this
pair admits a principled decomposition. Either $\reach{H_1} \neq
\reach{H_2}$, in which case the observation map $\obsmap$ has discarded
information that a less lossy observation map could in principle have
retained, and the failure is a failure of the particular observation
procedure; or $\reach{H_1} = \reach{H_2}$, in which case no observation
map factoring through reachable-set structure can distinguish $H_1$ from
$H_2$, and the failure is a failure of recoverability as such, not of
any particular reconstruction procedure.
\end{corollary}

Corollary \ref{cor:nonident-diagnosis} gives the precise form of the
question that motivates this entire research program's treatment of
reconstruction: whether a given case of reconstruction failure reflects
an inference that was performed incorrectly, or evidence that has
structurally ceased to exist. The corollary shows this is not merely a
useful distinction to draw informally but a genuine dichotomy, decidable
in principle by checking whether the colliding histories share a
reachable set.

\chapter{Conclusion}
\label{ch:conclusion}

The four theorems of this monograph were introduced, in
Chapter~\ref{ch:introduction}, as four questions one can ask of a single
admissibility relation. Having proven each in turn, it is worth stating
plainly what has and has not been established.

Representation Invariance (Theorem~\ref{thm:invariance}) established
that sameness of representation, properly understood, is sameness of
admissibility structure and nothing else; substrate is invisible to the
theorem except insofar as it constrains $R$. Admissibility Preservation
(Theorem~\ref{thm:preservation}) established that coherence of a
trajectory is not a separate property requiring its own definition, but
exactly the nondegeneracy of the admissibility relation examined along
that trajectory. Monotonic Relaxation (Theorem~\ref{thm:relaxation})
established that a scalar measure of inadmissibility decreases along
repair trajectories specifically, as a direct consequence of how repair
is already characterized elsewhere in this program, and Corollary
\ref{cor:relaxation-continuation} fixed the precise, non-automatic
relationship between this claim and the unrelated ordering relation of
monotonic continuation. Non-Identifiability
(Theorem~\ref{thm:nonident-structural}) established that failures of
historical reconstruction are not contingent misfortunes but a
structural guarantee whenever observation is constrained to reachable-set
structure and the space of possible histories exceeds what that
structure can encode, and Corollary~\ref{cor:nonident-diagnosis} gave a
principled way to diagnose, in any particular case, whether a
reconstruction failure reflects a correctable error or a genuine limit.

Two things follow from having stated these four theorems together rather
than separately, as they have variously appeared in earlier work in this
program. The first is economy: each theorem's proof drew only on
definitions already established elsewhere, and none required
introducing new primitive vocabulary. The second, and more important, is
that the four theorems are visibly not four independent facts about four
different subjects --- representation, coherence, repair, and
reconstruction --- but four consequences of one relation, $x \adm{R} y$,
examined respectively at a single state, along a trajectory, under a
restorative dynamics, and under a lossy projection. What varies across
the four chapters is not the object of study but the operation performed
on it. This is, in the end, the thesis this monograph exists to make
precise: that admissibility is not one tool among several available for
studying structure under transformation, but the common substrate from
which invariance, coherence, repair, and the limits of recovery are all
derivable in turn.


\backmatter

\begin{thebibliography}{9}

\bibitem{lamport1978}
Leslie Lamport, ``Time, Clocks, and the Ordering of Events in a
Distributed System,'' \textit{Communications of the ACM}, 1978.

\bibitem{chandy1985}
K.~Mani Chandy and Leslie Lamport, ``Distributed Snapshots: Determining
Global States of Distributed Systems,'' \textit{ACM Transactions on
Computer Systems}, 1985.

\bibitem{shannon1948}
Claude E.~Shannon, ``A Mathematical Theory of Communication,'' \textit{Bell
System Technical Journal}, 1948.

\bibitem{cover2006}
Thomas M.~Cover and Joy A.~Thomas, \textit{Elements of Information
Theory}, 2nd ed., Wiley, 2006.

\bibitem{friston2010}
Karl Friston, ``The Free-Energy Principle: A Unified Brain Theory?,''
\textit{Nature Reviews Neuroscience}, 2010.

\bibitem{aubin1991}
Jean-Pierre Aubin, \textit{Viability Theory}, Birkh\"auser, 1991.

\bibitem{powers1973}
William T.~Powers, \textit{Behavior: The Control of Perception}, Aldine,
1973.

\end{thebibliography}


\end{document}
